\chapter{Introduction}
\label{chap:intro}

% =====================================================================
\section{Context and Problem Statement}
\label{sec:context}
% =====================================================================

Understanding why things happen matters more than merely knowing that things happen together. When a doctor prescribes a treatment, an economist evaluates a policy, or an engineer diagnoses a system failure, the goal is not to find correlations but to understand causal mechanisms~\cite{pearl2009causality}. Knowing that treatment and recovery are correlated is insufficient for making good medical decisions; what matters is whether the treatment causes recovery or whether both are confounded by some third factor such as patient age or disease severity.

This thesis addresses the \textbf{ADIA Lab Causal Discovery Challenge}, a competitive machine learning benchmark organized by ADIA Lab in collaboration with CrunchDAO~\cite{crunchdao_causal_2024, crunchdao_adia_docs}. The task requires classifying the causal role of each variable in a dataset with a designated treatment $X$ and outcome $Y$. Given $n=1000$ observational samples over $p\in\{3,\ldots,10\}$ variables, the model must assign each remaining variable $K$ to one of eight causal roles: Confounder, Mediator, Collider, Independent, Cause of X, Consequence of X, Cause of Y, or Consequence of Y.

This eight-class classification problem presents four specific challenges that classical causal discovery methods struggle to overcome. First, subtle class distinctions make role pairs difficult to separate: Confounder ($K\to X, K\to Y$) and Cause of X ($K\to X$ only) both induce correlation between $K$ and $X$; similarly, Mediator ($X\to K\to Y$) and Consequence of X ($X\to K$ only) both exhibit dependence between $X$ and $K$. The distinguishing signal lies in subtle asymmetries that scalar statistics cannot capture. Second, class imbalance causes models to neglect minority roles such as Mediator and Collider without careful handling strategies. Third, permutation invariance requires the model to treat $(X,Y,K_1,K_2)$ identically to $(K_2,X,K_1,Y)$ since variable naming order carries no semantic information. Fourth, nonlinear relationships mean that simple Pearson correlation is often insufficient when causal roles differ in the shape rather than strength of their conditional relationships.

% =====================================================================
\section{Existing Solutions and Their Limitations}
\label{sec:limitations}
% =====================================================================

Prior work on this challenge falls into three groups, each with fundamental limitations that constrain performance.

\textbf{Classical Causal Discovery Methods} rely on conditional independence tests such as the PC algorithm to identify causal structure from observational data~\cite{hasan2023survey_arxiv}. These methods achieve only \textbf{~40\% Balanced Accuracy} on the ADIA Lab challenge~\cite{olivetti2025adialab}, barely above the 12.5\% random baseline for eight-class classification. Their limitations are severe: \textbf{limited to linear relationships and simple dependencies}, and \textbf{cannot handle the eight-class fine-grained classification task} since their output is a causal graph, not a role assignment for each variable.

\textbf{Machine Learning Feature Engineering Approaches} extract scalar statistics from each variable pair and train gradient boosted classifiers. A LightGBM model on 11 hand-crafted features achieves 63.11\%. A comprehensive pipeline with 300+ features including Pearson, Spearman, and distance correlation coefficients, partial correlations, conditional mutual information, HSIC, ANM scores, and outputs from PC, LiNGAM, NOTEARS, and GES algorithms reaches 72.64\% with XGBoost and LightGBM ensembles. Despite this extensive feature engineering, performance remains 8.44\% below the final model. The core problem is that \textbf{scalar features irreversibly discard curve shape information}: conditional expectation curves encode the functional form of causal relationships, and reducing them to single numbers loses the distinguishing patterns between roles.

\textbf{Deep Learning Approaches} began with the competition's first-place solution, which introduced 3-channel edge tensors encoding conditional distributions as structured inputs for a neural classifier~\cite{adia_rank1_report}. This approach achieves 76.70\% Balanced Accuracy, establishing a strong baseline that outperforms all classical and ML feature engineering methods. However, it has three critical limitations: \textbf{only 3 channels miss multi-scale patterns} that multi-bandwidth kernel regression can capture; \textbf{no structural awareness in the attention mechanism} means the model cannot encode topology-aware priors for chain, fork, and collider relationships; and \textbf{no 2D visual representation of joint distributions} leaves patterns invisible to pairwise analysis.

These limitations collectively define the performance ceiling of existing approaches. Addressing all eight requires a fundamentally different architecture that preserves curve shape information, captures multi-scale causal asymmetry, encodes structural inductive bias, and combines complementary visual and statistical representations.

% =====================================================================
\section{Our Solution and Results}
\label{sec:solution}
% =====================================================================

This thesis proposes a dual-pipeline neural architecture that addresses all eight limitations through five targeted contributions. The first pipeline processes 1D edge tensors encoding conditional distributions; the second processes 2D scatter density images encoding joint distribution topology. These pipelines are fused through learned gating and trained with a weighted dual-loss objective.

The final model achieves \textbf{81.082\% Balanced Accuracy} on the private test set, surpassing the first-place solution by +4.38 percentage points and outperforming the faithful baseline reimplementation by +7.12 percentage points.

% =====================================================================
\section{Five Technical Contributions}
\label{sec:contributions}
% =====================================================================

The first contribution is an \textbf{8-Channel Edge Tensor} that extends the baseline 3-channel representation by adding multi-bandwidth kernel regression channels ($h\in\{0.2,0.5,1.0\}$) capturing causal asymmetry at different scales, and ANM residual channels providing directional asymmetry signals. This representation captures patterns that a single bandwidth cannot, yielding +1.89\% improvement over the 3-channel baseline.

The second contribution is a \textbf{Structural Attention Bias} that injects topology-aware priors into self-attention layers. Rather than adding thousands of unconstrained parameters (which produced negative gains of -1.19\% to -1.90\% in failed experiments), this approach encodes six causal relationship types using only six learned scalars, outperforming those larger models by over 3.5 percentage points and providing the largest single architectural gain of +2.65\%.

The third contribution is a \textbf{2D Scatter Density Pipeline} that renders joint distributions as 8-channel $32\times32$ density images processed by a convolutional encoder. This complements the 1D edge tensor pipeline by capturing patterns such as multimodality and conditional heteroskedasticity that are invisible to pairwise scalar statistics, contributing +0.88\% with particularly strong gains on Confounder (+2.78\%) and Consequence of Y (+2.18\%).

The fourth contribution is \textbf{X/Y Remap Data Augmentation} that treats each directed edge $A\to B$ in the ground-truth DAG as a new treatment-outcome pair $(X',Y')$, recomputing all eight role labels accordingly. This multiplies effective training data by approximately 11x without additional data collection, yielding +1.55\% to +3.32\% depending on base model configuration.

The fifth contribution is a \textbf{Dual-Head Loss Weighting} scheme with $\lambda=0.7$ balancing the edge and node loss heads. This ensures both pipelines remain causally meaningful throughout training rather than one dominating, contributing an additional +0.59\% improvement.

\begin{table}[ht]
	\centering
	\caption{Performance comparison on the ADIA Lab Causal Discovery Challenge.
		Classical methods suffer from linear assumptions and coarse-grained output.
		ML feature engineering irreversibly loses curve shape when compressing to scalars.
		The 3-channel deep learning baseline misses multi-scale patterns, lacks structural awareness, and has no visual representation.
		The proposed method addresses all eight limitations simultaneously, achieving state-of-the-art performance.}
	\label{tab:sota_comparison}
	\begin{tabular}{lcc}
		\toprule
		\textbf{Approach}                                           & \textbf{Balanced Accuracy} & \textbf{Gap to SOTA} \\
		\midrule
		Classical baselines (PC, RF, DL)~\cite{olivetti2025adialab} & $\sim$40\%                 & $-$41\%              \\
		ML fullstack ensemble (300+ features)                       & 72.64\%                    & $-$8.44\%            \\
		Competition 1st place (3-channel)~\cite{adia_rank1_report}  & 76.70\%                    & $-$4.38\%            \\
		Proposed (8-channel + 2D + XY aug)                          & \textbf{81.082\%}          & \textbf{Baseline}    \\
		\bottomrule
	\end{tabular}
\end{table}

% =====================================================================
\section{Validation Evidence}
\label{sec:validation}
% =====================================================================

Ablation studies confirm that each contribution is necessary rather than redundant. Removing the edge context encoder collapses performance from 80.47\% to 51.51\%, a 29-percentage-point drop demonstrating that the 1D pipeline provides structurally necessary information that the 2D pipeline alone cannot capture. The four directional embeddings ($K\to X, K\to Y, X\to K, Y\to K$) propagate graph-wide causal context to each variable's representation.

The structural attention bias experiments provide a particularly instructive comparison. Adding 256 unconstrained attention parameters produced -1.19\% while adding 1024 parameters produced -1.90\%. By contrast, the six structured scalars encoding topology priors produced +2.65\%. This demonstrates that for this problem, structured inductive bias consistently outperforms raw model capacity when training data represents causal structure rather than arbitrary functions.

Per-class ablation analysis reveals complementary benefits across contributions. The 2D pipeline benefits Confounder (+2.78\%) and Mediator (+2.28\%) most, roles where joint distribution shape provides distinguishing information. The XY augmentation benefits rare classes most, validating that data multiplication addresses class imbalance rather than just overall volume.

% =====================================================================
\section{Thesis Organization}
\label{sec:organization}
% =====================================================================

The remainder of this thesis is organized as follows. Chapter~\ref{chap:background} reviews causal graph foundations, formally defines the eight causal roles, explains edge tensor construction and ANM residuals, and describes the deep learning components employed in this work. Chapter~\ref{chap:method} presents the complete dual-pipeline architecture including the 8-channel edge tensor, structural attention bias, 2D scatter density classifier, fusion mechanism, and training procedures. Chapter~\ref{chap:eval} describes the ADIA Lab dataset, evaluation metrics, baseline implementations, and experimental protocols. Chapter~\ref{chap:results} reports ablation studies, per-class performance analysis, mechanistic insights, and general lessons extracted from both successful and failed experiments.
